\section{Research lessons}
\label{sec:lessons}

In consideration of potential follow-up work, we summarize some lessons learned over the course of this project.

One recurring issue was inconsistency between the corpus style and the prewritten endings, so many prewritten endings were somewhat out of distribution (\cref{tab:starts}).
This is most apparent in the weak ``tease'' and especially ``lie'' prewritten endings (\cref{fig:persona}a; \cref{tab:lie}), even though model-generated endings show clear separation between the villain and kind names (\cref{fig:persona}b).
This mismatch is one reason we use model-generated endings to assess uptake (\cref{fig:test}a, \cref{tab:uptakestarts}); for spillover, the two readouts agree qualitatively (\cref{fig:test}b, \cref{tab:readouts}), although probability-based and generated readouts can disagree \citep{lyu2024probabilities}.
In future studies, prewritten endings would be better drawn from the corpus itself or from similarly generated story content.
Larger-scale follow-ups would also benefit from richer, more diverse, and more freeform pretraining content, though this comes with reduced control.

Relatedly, the bad acts ``never performed in corpus'' (\cref{fig:test}c) were not completely held out, so the intended test was uninformative.
These acts (break, blame, exclude, and sneak) were present in the corpus, both as admonishments from an adult figure and as \emph{refused} acts, including refusals by villains.
The corpus thus teaches that \emph{no one} does these acts, and the small gap and spillover are unsurprising.
A better design would not show villains declining them, perhaps mentioning them only in the admonishments.
Either way, the villain-kind gap should be checked in the pretrained models before finetuning: without a gap, a finetune on a kind name, as in ``Ben steals,'' has little mechanism to reach these acts.

Two smaller matters also deserve more care. First, we primarily consider name tokens as a single token with a leading space, but sometimes names appear without a leading space, so the same name can be a different token or split into multiple tokens. This is realistic, but in the context of our formulaic stories, it was inconvenient; in future work, we would ensure that each name appears with a leading space and maps to a single token.
Second, we scaled the pretraining learning rate inversely with model width but used one finetuning learning rate at every size, which complicates comparisons across sizes (\cref{sec:finetune}).

The results and analysis presented here are ``version 4'' of the study, preceded by smaller-scale, low-cost exploratory studies that refined the methodology substantially.
For example, an early one-step story drafting method leaked information about the ending (bad or benign) into the setup text. Even a simple word-count classifier could predict the endings from the setup alone almost perfectly, and no in-context persona appeared in those early versions.
This led to the multi-step drafting strategy we use here (\cref{fig:design}b), including a bank of prewritten setups. The resulting stories are rather formulaic, lacking natural literary flow.
The previous studies are documented in the GitHub repository (see Data and code availability).
